% Options for packages loaded elsewhere
% Options for packages loaded elsewhere
\PassOptionsToPackage{unicode}{hyperref}
\PassOptionsToPackage{hyphens}{url}
\PassOptionsToPackage{dvipsnames,svgnames,x11names}{xcolor}
%
\documentclass[
  english,
  11pt,
]{article}
\usepackage{xcolor}
\usepackage[margin=25mm]{geometry}
\usepackage{amsmath,amssymb}
\setcounter{secnumdepth}{-\maxdimen} % remove section numbering
\usepackage{iftex}
\ifPDFTeX
  \usepackage[T1]{fontenc}
  \usepackage[utf8]{inputenc}
  \usepackage{textcomp} % provide euro and other symbols
\else % if luatex or xetex
  \usepackage{unicode-math} % this also loads fontspec
  \defaultfontfeatures{Scale=MatchLowercase}
  \defaultfontfeatures[\rmfamily]{Ligatures=TeX,Scale=1}
\fi
\usepackage{lmodern}
\ifPDFTeX\else
  % xetex/luatex font selection
\fi
% Use upquote if available, for straight quotes in verbatim environments
\IfFileExists{upquote.sty}{\usepackage{upquote}}{}
\IfFileExists{microtype.sty}{% use microtype if available
  \usepackage[]{microtype}
  \UseMicrotypeSet[protrusion]{basicmath} % disable protrusion for tt fonts
}{}
\makeatletter
\@ifundefined{KOMAClassName}{% if non-KOMA class
  \IfFileExists{parskip.sty}{%
    \usepackage{parskip}
  }{% else
    \setlength{\parindent}{0pt}
    \setlength{\parskip}{6pt plus 2pt minus 1pt}}
}{% if KOMA class
  \KOMAoptions{parskip=half}}
\makeatother
% Make \paragraph and \subparagraph free-standing
\makeatletter
\ifx\paragraph\undefined\else
  \let\oldparagraph\paragraph
  \renewcommand{\paragraph}{
    \@ifstar
      \xxxParagraphStar
      \xxxParagraphNoStar
  }
  \newcommand{\xxxParagraphStar}[1]{\oldparagraph*{#1}\mbox{}}
  \newcommand{\xxxParagraphNoStar}[1]{\oldparagraph{#1}\mbox{}}
\fi
\ifx\subparagraph\undefined\else
  \let\oldsubparagraph\subparagraph
  \renewcommand{\subparagraph}{
    \@ifstar
      \xxxSubParagraphStar
      \xxxSubParagraphNoStar
  }
  \newcommand{\xxxSubParagraphStar}[1]{\oldsubparagraph*{#1}\mbox{}}
  \newcommand{\xxxSubParagraphNoStar}[1]{\oldsubparagraph{#1}\mbox{}}
\fi
\makeatother


\usepackage{longtable,booktabs,array}
\newcounter{none} % for unnumbered tables
\usepackage{calc} % for calculating minipage widths
% Correct order of tables after \paragraph or \subparagraph
\usepackage{etoolbox}
\makeatletter
\patchcmd\longtable{\par}{\if@noskipsec\mbox{}\fi\par}{}{}
\makeatother
% Allow footnotes in longtable head/foot
\IfFileExists{footnotehyper.sty}{\usepackage{footnotehyper}}{\usepackage{footnote}}
\makesavenoteenv{longtable}
\usepackage{graphicx}
\makeatletter
\newsavebox\pandoc@box
\newcommand*\pandocbounded[1]{% scales image to fit in text height/width
  \sbox\pandoc@box{#1}%
  \Gscale@div\@tempa{\textheight}{\dimexpr\ht\pandoc@box+\dp\pandoc@box\relax}%
  \Gscale@div\@tempb{\linewidth}{\wd\pandoc@box}%
  \ifdim\@tempb\p@<\@tempa\p@\let\@tempa\@tempb\fi% select the smaller of both
  \ifdim\@tempa\p@<\p@\scalebox{\@tempa}{\usebox\pandoc@box}%
  \else\usebox{\pandoc@box}%
  \fi%
}
% Set default figure placement to htbp
\def\fps@figure{htbp}
\makeatother



\ifLuaTeX
\usepackage[bidi=basic,shorthands=off]{babel}
\else
\usepackage[bidi=default,shorthands=off]{babel}
\fi
\ifLuaTeX
  \usepackage{selnolig} % disable illegal ligatures
\fi


\setlength{\emergencystretch}{3em} % prevent overfull lines

\providecommand{\tightlist}{%
  \setlength{\itemsep}{0pt}\setlength{\parskip}{0pt}}



 


\makeatletter
\@ifpackageloaded{caption}{}{\usepackage{caption}}
\AtBeginDocument{%
\ifdefined\contentsname
  \renewcommand*\contentsname{Table of contents}
\else
  \newcommand\contentsname{Table of contents}
\fi
\ifdefined\listfigurename
  \renewcommand*\listfigurename{List of Figures}
\else
  \newcommand\listfigurename{List of Figures}
\fi
\ifdefined\listtablename
  \renewcommand*\listtablename{List of Tables}
\else
  \newcommand\listtablename{List of Tables}
\fi
\ifdefined\figurename
  \renewcommand*\figurename{Figure}
\else
  \newcommand\figurename{Figure}
\fi
\ifdefined\tablename
  \renewcommand*\tablename{Table}
\else
  \newcommand\tablename{Table}
\fi
}
\@ifpackageloaded{float}{}{\usepackage{float}}
\floatstyle{ruled}
\@ifundefined{c@chapter}{\newfloat{codelisting}{h}{lop}}{\newfloat{codelisting}{h}{lop}[chapter]}
\floatname{codelisting}{Listing}
\newcommand*\listoflistings{\listof{codelisting}{List of Listings}}
\makeatother
\makeatletter
\makeatother
\makeatletter
\@ifpackageloaded{caption}{}{\usepackage{caption}}
\@ifpackageloaded{subcaption}{}{\usepackage{subcaption}}
\makeatother
\usepackage{bookmark}
\IfFileExists{xurl.sty}{\usepackage{xurl}}{} % add URL line breaks if available
\urlstyle{same}
% fallback for those not using the hyperref driver hyperxmp:
\makeatletter
\@ifundefined{xmpquote}{\newcommand{\xmpquote}[1]{#1}}{}
\makeatother
\hypersetup{
  pdftitle={How can motion with zero average transport momentum?},
  pdflang={en},
  colorlinks=true,
  linkcolor={blue},
  filecolor={Maroon},
  citecolor={Blue},
  urlcolor={Blue},
  pdfcreator={LaTeX via pandoc}}


\title{How can motion with zero average transport momentum?}
\usepackage{etoolbox}
\makeatletter
\providecommand{\subtitle}[1]{% add subtitle to \maketitle
  \apptocmd{\@title}{\par {\large #1 \par}}{}{}
}
\makeatother
\subtitle{Two opposite signs can produce the same useful effect}
\author{}
\date{2026-09-12}
\begin{document}
\maketitle


\href{../paper-map/article.html}{Start with the paper map} · Main route,
article 6 of 9

Imagine a road between two warehouses. Ten loaded trucks travel north
while ten empty trucks return south. The average direction of truck
travel can be zero, yet goods have clearly moved north.

Something similar happens in a fluid. Opposite motions need not cancel
the quantity they carry. Small-scale velocity variations can have zero
average and still produce a useful transport of momentum.

This is the essential reason the paper's oscillations can help its
background vortex. We will use signed numbers and multiplication to make
the idea precise.

\subsection{Velocity and transported cargo answer different
questions}\label{velocity-and-transported-cargo-answer-different-questions}

Averaging the truck velocities asks about their net movement. Counting
transported goods asks how much cargo crosses a boundary and in which
direction.

For a fluid, the relevant cargo is \textbf{momentum}: mass times
velocity. Fluid crossing an imaginary surface carries momentum with it.

This gives a product. One velocity component describes motion across the
surface. Another describes the momentum being carried along a chosen
direction. For constant density, their product gives the corresponding
momentum flux per unit density.

``Flux'' simply means an amount passing through a surface per unit area
and time. The direction of both the crossing and the transported
momentum matters.

\subsection{Try the arithmetic with two
fluctuations}\label{try-the-arithmetic-with-two-fluctuations}

Suppose equal portions of fluid have the following deviations from their
local mean motion:

{\def\LTcaptype{none} % do not increment counter
\begin{longtable}[]{@{}
  >{\raggedright\arraybackslash}p{(\linewidth - 6\tabcolsep) * \real{0.2500}}
  >{\raggedright\arraybackslash}p{(\linewidth - 6\tabcolsep) * \real{0.2500}}
  >{\raggedright\arraybackslash}p{(\linewidth - 6\tabcolsep) * \real{0.2500}}
  >{\raggedright\arraybackslash}p{(\linewidth - 6\tabcolsep) * \real{0.2500}}@{}}
\toprule\noalign{}
\begin{minipage}[b]{\linewidth}\raggedright
Portion
\end{minipage} & \begin{minipage}[b]{\linewidth}\raggedright
Upward velocity fluctuation
\end{minipage} & \begin{minipage}[b]{\linewidth}\raggedright
Horizontal velocity fluctuation
\end{minipage} & \begin{minipage}[b]{\linewidth}\raggedright
Product
\end{minipage} \\
\midrule\noalign{}
\endhead
\bottomrule\noalign{}
\endlastfoot
A & +2 & +3 & +6 \\
B & −2 & −3 & +6 \\
Average & 0 & 0 & +6 \\
\end{longtable}
}

Both velocity fluctuations average to zero. Their product does not.

For portion A, upward motion carries an excess of horizontal momentum.
For portion B, downward motion carries a deficit. Relative to the mean
flow, these contribute to the same signed transport.

The arithmetic is simply

\[
(+2)(+3)=6,\qquad (-2)(-3)=6.
\]

\begin{figure}

\centering{

\pandocbounded{\includegraphics[keepaspectratio,alt={Two rows show fluctuations plus two and plus three, then minus two and minus three; both products are plus six.}]{figures/reader.png}}

}

\caption{\label{fig-transport}Opposite fluctuations have zero mean,
while the products in the transport column have the same sign. These
invented values illustrate the averaging mechanism.}

\end{figure}%

A sine wave makes the same point smoothly: its positive and negative
parts cancel in the mean, while its square is never negative. The
supporting notebook uses sine waves to construct and check this kind of
transport.

\subsection{Why this averaged transport is called a
stress}\label{why-this-averaged-transport-is-called-a-stress}

Different surfaces and momentum directions give different products.
Collecting them describes how small-scale motion transports momentum
through the fluid.

That collection is often called a \textbf{Reynolds stress}, with sign
and density conventions that vary between texts. For this explanation,
keep the physical meaning: it is the transport contributed by the
velocity fluctuations.

You do not need to learn matrix notation to understand its job. The
paper especially needs radial transport of circular momentum and radial
transport of axial momentum.

The word ``radial'' means toward or away from the axis. ``Axial'' means
along the axis. Those are two different kinds of cargo carried across
rings in the vortex.

\subsection{Transport must change from place to place to create a net
effect}\label{transport-must-change-from-place-to-place-to-create-a-net-effect}

Imagine a warehouse receiving 100 boxes an hour and sending out 100. Its
stock does not change. If it receives 100 and sends out 90, ten boxes
accumulate each hour.

Momentum obeys the same accounting. If equal momentum enters and leaves
a region, the transport does not change its total. A difference between
incoming and outgoing transport produces a net effect.

This is why a nonzero \textbf{constant} averaged stress does not by
itself accelerate the mean flow. The useful force arises from how the
transport changes across space.

The paper shapes the pulse amplitudes and locations so their transport
has the required spatial variation. It aims to cancel the troublesome
contribution identified by the
\href{../d3_residual/article.html}{residual calculation}.
\href{https://cdn.openai.com/pdf/32d9f210-8b73-45e0-91bc-82a30aef8a9a/navier-stokes.pdf}{Paper,
Section 2.2 and Section 3.3}

\subsection{Tune the amount with the wave
amplitude}\label{tune-the-amount-with-the-wave-amplitude}

If both velocity fluctuations are multiplied by two, their product is
multiplied by four. Therefore, changing the wave amplitude changes its
quadratic transport strength.

A desired contribution four times as large needs twice the amplitude in
this simple scaling. That square-root relationship explains why the
paper first finds positive weights for the required stress and then uses
their square roots as amplitude factors.

The two pulse families provide different ratios of circular to axial
momentum transport. Their nonnegative weights combine those
contributions. The desired pair has to be attainable with those
families---the ``stress cone'' condition explained in the
\href{../convex-integration/article.html}{convex-integration article}.

The paper does not claim that any arbitrary ripple can supply any
arbitrary force. Directions, weights, geometry, and evolution must all
fit.

\subsection{Keep the fluid-volume constraint
intact}\label{keep-the-fluid-volume-constraint-intact}

Simply making a wave stronger in one place can accidentally introduce an
apparent accumulation or loss of fluid volume.

The source therefore builds its corrections through a representation
that guarantees incompressibility. This introduces some additional,
smaller velocity terms. They are retained in the next residual
calculation rather than ignored.

Our numerical examples use a related construction to check zero mean,
the desired averaged products, and the effect of a spatially changing
amplitude. They illustrate the transport mechanism in a simple periodic
domain. They are not the actual pulse field of the full vortex.

\subsection{Where this fits in the bigger
picture}\label{where-this-fits-in-the-bigger-picture}

\textbf{The oscillations can supply a missing mean transport even though
their mean velocity is zero.} This connects fine-scale motion to the
large-scale force balance.

That is a central bridge between the convex-integration viewpoint and
the paper's pulse construction. But the waves also need a source of
energy and a controlled way to fade. Continue with
\href{../d5_shearing_wave/article.html}{how a small wave grows and then
dies away}.

The \href{technical-article.html}{technical article},
\href{../../code/d4_mean_stress/study.ipynb}{notebook}, and
\href{source-notes.md}{source notes} contain the explicit waves,
averaging conventions, and computational checks.




\end{document}
